\section{Operator Overloading}

\subsection{Section Overview}
Operator overloading allows us to define how C++ operators behave with user-defined types (classes). This can make our custom classes more intuitive and easier to use, allowing them to be manipulated with familiar syntax like \texttt{+}, \texttt{-}, \texttt{==}, \texttt{<<}, etc. In this section, we'll learn the rules and syntax for overloading various operators.

\subsection{What is Operator Overloading?}
Operator overloading is a form of polymorphism where we provide a special meaning for an operator for a given user-defined type. It allows us to integrate our custom types seamlessly with the language. For example, we can overload the \texttt{+} operator to add two \texttt{Vector} objects or the \texttt{<<} operator to print a \texttt{Player} object to the console.

\begin{remark}
	Not every operator can be overloaded. The following operators \textbf{cannot} be overloaded: \texttt{::} (scope resolution), \texttt{.} (member access), \texttt{.*} (pointer-to-member), and \texttt{?:} (ternary conditional). Also, you cannot create new operators, change the arity of an operator, or change its precedence.
\end{remark}

\subsection{Overloading the Assignment Operator (copy)}
The copy assignment operator (\texttt{operator=}) is used to copy the values from one object to another already existing object. If your class manages dynamic memory (raw pointers), you must provide a deep copy implementation to avoid shallow copy issues.
\begin{lstlisting}
	class MyString {
	private:
		char *str;
	public:
		// ... other members ...
		MyString& operator=(const MyString &rhs); // Copy assignment
	};

	MyString& MyString::operator=(const MyString &rhs) {
		if (this == &rhs) { // self-assignment check
			return *this;
		}
		delete[] str; // deallocate old storage
		str = new char[strlen(rhs.str) + 1]; // allocate new storage
		strcpy(str, rhs.str); // copy
		return *this; // return reference to this object
	}
\end{lstlisting}

\subsection{Overloading the Assignment Operator (move)}
C++11 introduced the move assignment operator (\texttt{operator=}). It's used to transfer ownership of resources from a temporary (r-value) object to another object, avoiding a costly deep copy.
\begin{lstlisting}
	class MyString {
	public:
		// ... other members ...
		MyString& operator=(MyString &&rhs); // Move assignment
	};

	MyString& MyString::operator=(MyString &&rhs) {
		if (this == &rhs) { // self-assignment check
			return *this;
		}
		delete[] str; // deallocate old storage
		str = rhs.str; // steal the pointer
		rhs.str = nullptr; // null out the source pointer
		return *this;
	}
\end{lstlisting}

\subsection{Overloading Operators as Member Functions}
Many operators are best overloaded as member functions. The left-hand operand is the object the method is called on (\texttt{*this}). The right-hand operand is passed as a parameter.
\begin{lstlisting}
	class Money {
	public:
		int dollars;
		int cents;

		Money operator+(const Money &rhs) const {
			int total_cents = (this->dollars * 100 + this->cents)
			                + (rhs.dollars * 100 + rhs.cents);
			return Money{total_cents / 100, total_cents % 100};
		}
	};

	Money m1 {1, 50}, m2 {2, 75};
	Money m3 = m1 + m2; // m1.operator+(m2)
\end{lstlisting}

\subsection{Overloading Operators as Global Functions}
Some operators must be overloaded as global (non-member) functions. This is common when the left-hand operand is not a class object, or when you want automatic type conversion on the left-hand operand. Stream insertion/extraction operators are a prime example.
\begin{lstlisting}
	class Money { /* ... */ };

	Money operator+(const Money &lhs, const Money &rhs) {
		// ...
	}
\end{lstlisting}

\subsection{Coding Exercise 35: Overloading the Addition Operator}
Given the \texttt{Account} class, overload the \texttt{+} operator as a member function to add two \texttt{Account} objects, returning a new \texttt{Account} with the combined balance.
\begin{lstlisting}
	#include <iostream>

	class Account {
	public:
		double balance;
		Account(double b = 0.0) : balance(b) {}

		Account operator+(const Account &rhs) const {
			return Account(this->balance + rhs.balance);
		}
	};

	int main() {
		Account a1(100.0);
		Account a2(250.50);
		Account a3 = a1 + a2;
		std::cout << "Combined balance: " << a3.balance << std::endl; // 350.50
		return 0;
	}
\end{lstlisting}

\subsection{Overloading Stream Insertion and Extraction Operators}
To allow your custom objects to work with \texttt{std::cout} and \texttt{std::cin}, you should overload the \texttt{<<} and \texttt{>>} operators. These must be overloaded as global functions.
\begin{lstlisting}
	class Money {
	public:
		int dollars;
		int cents;
		// ... constructor, etc ...
		friend std::ostream& operator<<(std::ostream &os, const Money &money);
		friend std::istream& operator>>(std::istream &is, Money &money);
	};

	std::ostream& operator<<(std::ostream &os, const Money &money) {
		os << "$" << money.dollars << "." << money.cents;
		return os;
	}

	std::istream& operator>>(std::istream &is, Money &money) {
		is >> money.dollars >> money.cents;
		return is;
	}

	int main() {
		Money m{10, 50};
		std::cout << "I have " << m << std::endl; // I have $10.50
	}
\end{lstlisting}

\subsection{Section Challenge}
Starting with the \texttt{MyString} class, overload the following operators:
\begin{itemize}
	\item \texttt{-} (unary minus): returns a lowercase version of the string.
	\item \texttt{==} (equality): returns true if two \texttt{MyString} objects are the same.
	\item \texttt{+} (concatenation): returns a new \texttt{MyString} object that is the concatenation of two others.
\end{itemize}

\subsection{Section Challenge --- Solution}
\begin{lstlisting}
	#include <iostream>
	#include <cstring>
	#include <cctype>

	// NOTE: This class intentionally omits the copy constructor and copy
	// assignment operator to keep the example focused. A production-ready
	// version would need all three (Rule of Three) to avoid shallow-copy
	// double-free bugs when copying MyString objects.
	class MyString {
	private:
		char *str;
	public:
		MyString(const char *s = nullptr) {
			if (s) {
				str = new char[strlen(s) + 1];
				strcpy(str, s);
			} else {
				str = new char[1];
				*str = '\0';
			}
		}
		~MyString() { delete[] str; }

		void display() const { std::cout << str; }

		// Unary minus - lowercase
		MyString operator-() const {
			char *buff = new char[strlen(str) + 1];
			strcpy(buff, str);
			size_t len = strlen(buff); // compute once, not every iteration
			for (size_t i = 0; i < len; ++i) {
				buff[i] = tolower(buff[i]);
			}
			MyString temp {buff};
			delete[] buff;
			return temp;
		}

		// Equality
		bool operator==(const MyString &rhs) const {
			return (strcmp(this->str, rhs.str) == 0);
		}

		// Concatenation
		MyString operator+(const MyString &rhs) const {
			char *buff = new char[strlen(this->str) + strlen(rhs.str) + 1];
			strcpy(buff, this->str);
			strcat(buff, rhs.str);
			MyString temp {buff};
			delete[] buff;
			return temp;
		}
	};

	int main() {
		MyString s1 {"HELLO"};
		MyString s2 = -s1; // lowercase
		s1.display(); std::cout << std::endl; // HELLO
		s2.display(); std::cout << std::endl; // hello

		MyString s3 {"world"};
		MyString s4 = s2 + " " + s3;
		s4.display(); std::cout << std::endl; // hello world

		std::cout << std::boolalpha;
		std::cout << (s1 == "HELLO") << std::endl; // true
		std::cout << (s2 == s3) << std::endl; // false

		return 0;
	}
\end{lstlisting}

\subsection{Quiz 11: Section 14 Quiz}
\begin{enumerate}
	\item Why would you overload an operator for a custom class?
	\begin{enumerate}
		\item[a)] To make the code harder to read.
		\item[b)] It's required by the compiler for all classes.
		\item[c)] To provide intuitive, natural syntax for using objects of the class.
		\item[d)] To change the behavior of operators for built-in types like \texttt{int}.
	\end{enumerate}

	\item Which operators must be overloaded as global (non-member) functions?
	\begin{enumerate}
		\item[a)] All unary operators.
		\item[b)] The assignment operator \texttt{=}.
		\item[c)] The stream insertion \texttt{<<} and extraction \texttt{>>} operators.
		\item[d)] No operators have to be global.
	\end{enumerate}

	\item When overloading the copy assignment \texttt{operator=}, what is an important first step in the implementation?
	\begin{enumerate}
		\item[a)] Allocating new memory.
		\item[b)] Checking for self-assignment.
		\item[c)] Returning \texttt{*this}.
		\item[d)] Deleting the incoming object.
	\end{enumerate}
\end{enumerate}

\textbf{Answers:} 1. (c), 2. (c), 3. (b)
